%------------------------------------------------------------------------------%
\section{Exercícios}

\begin{enumerate}

\item Determine a equação da reta que passa pelos pontos dados:
\begin{itemize}
\item[a)] $(4,1)$ e $(2,3)$
\item[b)] $(1,2)$ e $(-2,-4)$
\end{itemize}

\item Um clube vem perdendo sócios a um ritmo constante ano após ano, tendo registrado 14 mil sócios em 2004 e apenas 8 mil sócios em 2013.

\begin{itemize}
\item[a)] Defina uma função linear (ou afim) $S(t)$ que forneça o número de sócios do clube no ano $t$, supondo que $t$ seja zero no ano 2000.
\item[b)] O clube pretende fechar as portas se o número de sócios chegar a 2mil. Usando a função encontrada no item a), determine em que ano isso ocorrerá.

\item[c)] Construa o gráfico de $S$.
\end{itemize}
\item Sejam dadas as funções $f(x)= \frac{1}{6}x-2$ e $g(x)=3-\frac{2}{3}x$.
\begin{itemize}
\item[a)] Exiba os gráficos de $f$ e $g$ no mesmo plano cartesiano.
\item[b)] Determine para que valores de $x$ a desigualdade $f(x) \leq g(x)$ é verdadeira.
\end{itemize}
\item (Minimização da área cercada) Um fazendeiro pretende usar $500m$ de cerca para proteger um bosque retagular de comprimento $x$ e largura $y$, às margens de um riacho.

\begin{itemize}
\item[a)] Usando o comprimento da cerca, escreva o valor de $y$ em função de $x$.
\item[b)] Com base na expressão encontrada no item a), escreva a função $A(x)$ que fornece a área cercada com relação a $x$.
\item[c)] Determine o valor de $x$ que maximiza a área cercada. Determine também o valor de $y$ e a área máxima.
\item[d] Trace o gráfico de $A(x)$.
\end{itemize}

\item Dada a função $f(x)= x^2-3x$,
\begin{itemize}
\item[a)] Determine algebricamente os pontos no quais $f(x)=0$.
\item[b)]Determine algebricamente os pontos no quais $f(x)=-2$.
\item[c)] Esboce o gráfico da função no plano coordenado, indicando os pontos que foram obtidos no item anterior.
\item[d] Determine graficamente a solução da inequação $f(x) \geq -2$. Confirme sua resposta resolvendo analiticamente a inequação.

\end{itemize}

\item Dada a função $f(x)= 15x^2+x-2$,
\begin{itemize}
\item[a)] Determine algebricamente os pontos no quais $f(x)=0$.
\item[b)]Determine algebricamente os pontos no quais $f(x) \leq-2$.
\item[c)] determine algebricamente os pontos de mínimo ou máximo de $f$.
\item[d)] Esboce o gráfico de $f$

\end{itemize}

\item O lucro(em milhões de reais) que uma fábrica obtém com a venda de um produto é dado pela função $L(x)=-\frac{x^2}{2}+3x+6$, em que $x$ é o valor gasto (também em milhões de reais) com propaganda de televisão. Calcule o valor que a empresa deve gastar com propaganda para obter o lucro máximo. Determine o lucro nesse caso.


\item Faça o esboço da parábolas a seguir.

\begin{minipage}[t]{.44\textwidth}
\begin{itemize}
\item[a)] $ f(x)=x^2-3x+4$
\item[b)] $f(x)=-2x^2+3x+2$
\item[c)] $f(x)=2x+x^2$
\end{itemize}
\end{minipage}
\begin{minipage}[t]{.44\textwidth}
\begin{itemize}
\item[d)]$f(x)= -x^2-4$
\item[e)]  $f(x)=(x-4)(x+1)$
\item[f)] $f(x)=-2x^2+8$
\end{itemize}
\end{minipage}

 \item Deve-se construir um tanque de aço em forma de um cilindro circular reto de 3m de altura com
dois hemisférios nos extremos. O raio r ainda está por determinar. Expresse a área S da
superfície do tanque em função de r.
\item Desenvolva as expressões:

\begin{minipage}[t]{.44\textwidth}
\begin{itemize}
\item[a)] $ (x-3)^2$
\item[b)] $ (\sqrt{x}+2)^2$
\item[c)] $(x^2-\sqrt{2})(x^2+\sqrt{2})$
\item[d)] $(x-1)(x^2-5)$
\item[e)] $(\sqrt{3}-\sqrt{2})(\sqrt{3}-5)$
\end{itemize}
\end{minipage}
\begin{minipage}[t]{.44\textwidth}
\begin{itemize}
\item[f)] $(\sqrt{3}-\sqrt{2})(\sqrt{3}+\sqrt{2})$
\item[g)] $(3x-2)^2$
\item[h)] $ (\sqrt{x}-\sqrt{3})^2$
\item[i)] $x(2xy-5)(x-y)$
\item[j)] $(x-4y)^3$
\end{itemize}
\end{minipage}
\item Simplifique as expressões, fatorando os termos, caso necessário. Suponha sempre que os denominadores são não nulos.

\begin{minipage}[t]{.44\textwidth}
\begin{itemize}
\item[a)] $ \frac{2x-6}{x-3}$
\item[b)] $ \frac{x^2-x^3}{x}$
\item[c)] $\frac{x^2-3x}{4x-12}$
\item[d)] $\frac{2x^2-50}{x^3+5x^2}$
\item[e)] $\frac{x^3-8}{x-2}$
\end{itemize}
\end{minipage}
\begin{minipage}[t]{.44\textwidth}
\begin{itemize}
\item[f)] $ \frac{x^2y-xy^2}{xy}$
\item[g)] $ \frac{x^2-9}{x^2-3x}$
\item[h)] $\frac{x^4-y^4}{x-y}$
\item[i)] $\frac{x^2-5x+4}{x-4}$
\item[j)] $\frac{x-9}{\sqrt{x}+3}$
\end{itemize}
\end{minipage}

\item Encontre as raízes dos polinômios a seguir e escreva-os na forma fatorada.

\begin{minipage}[t]{.44\textwidth}
\begin{itemize}
\item[a)] $ p(x)=9x^4-10x^2+1$
\item[b)] $ q(x)=x^3-4x^2-17x+60$
\end{itemize}
\end{minipage}
\begin{minipage}[t]{.44\textwidth}
\begin{itemize}
\item[c)] $g(x)=2x^3-3x^2-23x+12$
\item[d)] $ f(x)=x^4-4x^3+5x^2-2x$
\end{itemize}
\end{minipage}

\item Encontre o resto e o quociente da divisão dos seguintes polinômios

\begin{minipage}[t]{.44\textwidth}
\begin{itemize}
\item[a)] $ p(x)=2x^3-3x^2+6$ por $d(x)=2x-4$
\item[b)]  $ p(x)=6x^4+5x^3-2x$ por $d(x)=3x-2$
\end{itemize}
\end{minipage}
\begin{minipage}[t]{.44\textwidth}
\begin{itemize}
\item[c)] $g(x)=6x^2+7x+9$ por $2x^2-5x+1$
\end{itemize}
\end{minipage}

\item Resolva as inequações

\begin{minipage}[t]{.44\textwidth}
\begin{itemize}
\item[a)] $ 3x-5 \leq 0$
\item[b)]  $ (x+2)(x-3) \leq 0$
\item[c)] $x^2-3x \geq 0$
\end{itemize}
\end{minipage}
\begin{minipage}[t]{.44\textwidth}
\begin{itemize}
\item[d)] $-4x^2+12x-9 \leq 0$
\item[e)]$ \frac{x-2}{x+3} \leq 0$
\item[f)] $\frac{x+8}{x^2+7x+12} \leq 0$
\end{itemize}
\end{minipage}
\end{enumerate}

%------------------------------------------------------------------------------%
\subsection{Gabarito:}

\begin{enumerate}

\item
    \begin{enumerate}
        \item $y=-x+5$
        \item $y=2x$
    \end{enumerate}

\item
    \begin{enumerate}
        \item $S(t)=-\frac{2}{3}t+ \frac{50}{3}$
        \item $2022$
        \item $ $
            \begin{figure}[h]
                \centering
                \includegraphics[width=0.6\textwidth]{gabarito_aula2-2c}
                \caption{Gráfico da função $S(t)=-\frac{2}{3}t+ \frac{50}{3}$ }
            \end{figure}
    \end{enumerate}

\item
    \begin{enumerate}
        \item $ $
            \begin{figure}[h]
                \centering
                \includegraphics[width=0.6\textwidth]{gabarito_aula2-3a}
                \caption{Gráfico das funções $f(x)$ e $g(x)$ }
            \end{figure}
        \item $x \leq 6$
    \end{enumerate}

\item
    \begin{enumerate}
        \item $y=250-x$
        \item $A(x)=250x-x^2$
        \item $x=y=125m$ e $A_{max}=15625m^2$
        \item $ $
            \begin{figure}[h]
                \centering
                \includegraphics[width=0.6\textwidth]{gabarito_aula2-4d}
                \caption{Gráfico da função $A(x)$ }
            \end{figure}
    \end{enumerate}


\item
    \begin{enumerate}
        \item $x=0$ e $x=3$
        \item $x=1$ e $x=2$
        \item $ $
            \begin{figure}[h]
                \centering
                \includegraphics[width=0.6\textwidth]{gabarito_aula2-5c}
                \caption{Gráfico da função $f(x) = x^2-3x$ }
            \end{figure}
        \item $x \leq 1$ ou $x \geq 2$
    \end{enumerate}

\item
    \begin{enumerate}
        \item $x=-\frac{2}{5}$ e $x=\frac{1}{3}$
        \item $[-\frac{1}{15},0]$
        \item Ponto de mínimo: $(-\frac{1}{30},-\frac{121}{60})$
                \item $ $
            \begin{figure}[h]
                \centering
                \includegraphics[width=0.6\textwidth]{gabarito_aula2-6d}
                \caption{Gráfico da função $f(x) = 15x^2+x-1$ }
            \end{figure}
    \end{enumerate}

\item Gasto: $3$ milhões de reais \\ Lucro Máximo: $10,5$ milhões de reais

\item
    \begin{enumerate}
        \item
            \begin{figure}[h]
                \centering
                \includegraphics[width=0.6\textwidth]{exparaba}
                \caption{Gráfico da função $f(x) = x^2-3x+4$ }
            \end{figure}

        \item
            \begin{figure}[h]
                \centering
                \includegraphics[width=0.6\textwidth]{exparabb}
                \caption{Gráfico da função $f(x) = -2x^2+3x+2$ }
            \end{figure}
        \item
            \begin{figure}[h]
                \centering
                \includegraphics[width=0.6\textwidth]{exparabc}
                \caption{Gráfico da função $f(x) = 2x+x^2$ }
            \end{figure}
        \item
            \begin{figure}[h]
                \centering
                \includegraphics[width=0.6\textwidth]{exparabd}
                \caption{Gráfico da função $f(x) = -x^2-4$ }
            \end{figure}

        \item $ $
            \begin{figure}[h]
                \centering
                \includegraphics[width=0.6\textwidth]{exparabe}
                \caption{Gráfico da função $f(x) = (x-4)(x+1)$ }
            \end{figure}
        \item $ $
            \begin{figure}[h]
                \centering
                \includegraphics[width=0.6\textwidth]{exparabf}
                \caption{Gráfico da função $f(x) = -2x^2+8$ }
            \end{figure}
    \end{enumerate}

\item $S(r)=4\pi r^2+6\pi r$

\item
    \begin{enumerate}
        \item $x^2-6x+9$
        \item $x+4x+4$
        \item $x^4-2$
        \item $x^3-x^2-5x+5$
        \item $3-5\sqrt{3}-\sqrt{6}+5\sqrt{2}$
        \item $1$
        \item $9x^2-12x+4$
        \item $x-2\sqrt{3x}+3$
        \item $2x^3y-2x^2y^2-5x^2+5xy$
        \item $x^3-12x^2y+48xy^2-64y^3$
    \end{enumerate}

\item
    \begin{enumerate}
        \item $2$, se $x \neq 3$
        \item $x(1-x)$, se $x \neq 0$
        \item $\frac{x}{4}$, se $x \neq 3$
        \item $\frac{2(x-5)}{x^2}$, se $x \neq 5$
        \item $x^2+2x+4$, se $x \neq 2$
        \item $x-y$, se $xy \neq 0$
        \item $\frac{x+3}{x}$, se $x \neq 3$
        \item $(x+y)(x^2+y^2)$, se $x-y \neq 0$
        \item $x-1$, se $x \neq 4$
        \item $\sqrt{x}-3$
    \end{enumerate}

\item
    \begin{enumerate}
        \item Raízes: $x=\pm 1, x=\pm \frac{1}{3}$. Fatoração: $p(x)=9(x+1)(x-1)(x-\frac{1}{3})(x+\frac{1}{3})$
        \item Raízes: $x=-4, 3, 5$. Fatoração: $q(x)=(x+4)(x-3)(x-5)$
        \item Raízes: $x=-3, \frac{1}{2}, 4$. Fatoração: $p(x)=2(x+3)(x-\frac{1}{2})(x-4)$
        \item Raízes: $x=0, 1, 2$. Fatoração: $p(x)=x(x-1)^2(x-2)$
    \end{enumerate}

\item
    \begin{enumerate}
        \item $Q(x)=x^2+ \frac{1}{2}x+2, \,\, R(x)=10$
        \item $Q(x)=2x^3+ 3x^2+2x+\frac{2}{3}, \,\, R(x)=\frac{4}{3}$
        \item $Q(x)=3, \,\, R(x)=22x+6$
    \end{enumerate}

\item
    \begin{enumerate}
        \item $S=\{ x \in \R| x \leq \frac{5}{3} \}$
        \item $S=[-2,3]$
        \item $S=\{ x \in \R| x \leq 0 \,\, \text{ou} \,\, x \geq 3 \}$
        \item $S= \R$
        \item $S=\{ x \in \R|-3< x \leq 2 \}$
        \item $S=\{ x \in \R| x \leq -8 \,\, \text{ou} \,\, -4<x<-3 \}$
    \end{enumerate}

\end{enumerate}
